\documentclass[11pt]{article}
\usepackage[a4paper,margin=21mm]{geometry}
\usepackage[no-math]{fontspec}
\setmainfont{Latin Modern Roman}
\newfontfamily\CJKfont{Noto Serif CJK SC}
\XeTeXlinebreaklocale "zh"
\XeTeXlinebreakskip=0pt plus 1pt
\usepackage{amsmath,amssymb,amsthm,amscd,graphicx,color}
\usepackage{xurl}
\usepackage[unicode,hidelinks]{hyperref}
\allowdisplaybreaks
\setlength\emergencystretch{3em}
\numberwithin{equation}{section}

\DeclareMathOperator{\sgn}{sgn}
\newcommand\abs[1]{\lvert #1 \rvert}

\numberwithin{equation}{section}

\newtheorem{Theorem}{Theorem}[section]
\newtheorem{Corollary}[Theorem]{Corollary}
\newtheorem{Lemma}[Theorem]{Lemma}
 { \theoremstyle{definition}
\newtheorem{Remark}[Theorem]{Remark} }

\newfontfamily\nativebbone{DejaVu Math TeX Gyre}
\ExplSyntaxOn
\NewDocumentCommand{\mathbbm}{m}{\str_if_eq:nnTF{#1}{1}{\mathord{\text{\upshape\nativebbone\char"1D7D9}}}{\PackageError{nativecompat}{Unsupported~mathbbm~argument}{Only~verified~digit~one~is~accepted}}}
\ExplSyntaxOff

\makeatletter\newlabel{eq:con}{{2.3}{}{Original source locator}{}{}}\makeatother
\begin{document}
\begin{center}\Large The Heisenberg-group two-body Hamiltonian with nonzero coupling and positive masses is not rationally Liouville-integrable on its algebraic phase cover\end{center}
\noindent\textbf{Stable ID:} 1430\quad\textbf{Author:} Tomasz Stachowiak; Andrzej J. Maciejewski\par
\noindent\textbf{Work:} Non-Integrability of the Kepler and the Two-Body Problems on the Heisenberg Group\par
\noindent\textbf{Source:} \url{https://sigma-journal.com/2021/074/}\par
\noindent\textbf{Version:} Official SIGMA 17 (2021) article 074, DOI 10.3842/SIGMA.2021.074; journal PDF and native TeX acquired 2026-09-30, exact bytes pinned by SHA256\par
\noindent\textbf{Rights:} CC BY 4.0. Authors retain copyright. \url{https://creativecommons.org/licenses/by/4.0/}. Reuse and typesetting adaptation; original language and mathematical content preserved.\par
\noindent\textbf{Difficulty (prior selection preserved):} {\CJKfont H2: The proof chooses a nonconstant complex trajectory and explicitly reduces its four-dimensional normal variational subsystem by a time-dependent triangularizing matrix. An exponential gauge then produces a parabolic-cylinder equation with non-solvable differential Galois identity component, yielding the integrability obstruction. The algebraic potential requires the supplied Poisson phase-space extension and the precise rational-integral limitation at irregular infinity.}\par
\medskip\noindent\textbf{Editorial boundary note (not author text):} Selected stable problem is the positive-mass/nonzero real coupling rational-integral scope of printed Theorem3.1, on the nonsingular algebraic phase cover of AppendixA. Original theorem/proof language remains literal, while the author non-Fuchsian rational-only limitation2.1 is retained explicitly; no arbitrary local-meromorphic or all-energy-level obstruction is added. All group/rho/Hamiltonian/canonical variable/trajectory/rescaling and every variational/Q/subsystem/exponential-gauge/parabolic-cylinder parameter formula are included. Standard Morales–Ramis/Rehm/dAlambert prerequisites stay precise author references. Complete explicit algebraic cover/rational K/Poisson matrix/rank/Casimir/correspondence statement is retained, including the literal native729 direct-proof omission and prior Section2 pointer; root classifies that exact P=0 projection as a bounded finite implicit-differentiation/chain-rule prerequisite, not a separately supplied appendix proof. Its precise prior bibliographic identity/Section2 pointer and author omission remain prominent, without an editor verification. The mu=-1 Maple-only branch and one-body/factorization results are unselected. All three identity-matrix mathbbm\{1\} uses are preserved, with the previously authorized exact digit-one-only named DejaVu Math TeX Gyre U+1D7D9 typography resolver, rejecting all other arguments; font dependency and original glyph comparison disclosed. No mathematical formula changes.\par
\noindent\textbf{External original reference eq:con:} Original unselected one-body axial/radial potential condition native177–182: partial\_z V(0,0,c)=2a nonzero. It is mentioned only in the unselected open-potential contextual examples, not an extra selected premise.\par
\medskip\hrule\medskip\noindent\textbf{Author text begins below.}\par

% Original source lines 55--116
\section{Introduction}
The idea of studying the Kepler problem, or the problem of $n$ bodies, in a non-Euclidean space has a long history, but in almost all cases such
generalisations were performed in spaces of~con\-stant curvature. For a very detailed and critical recent overview
of this subject we refer to the nice article \cite{borisov:16::}.

In \cite{1} the authors considered the question of generalization of
the Kepler problem and its resulting mechanics, based on first principles. The idea was to
formulate an analog of the classical problem in spaces which are
homogeneous, isotropic and admit dilations. The last requirement is
very restrictive because among homogeneous Riemannian manifolds only
Euclidean ones admit dilations. The simplest non-Euclidean metric
space satisfying all the required properties is the Heisenberg group
of special upper-triangular matrices
\begin{equation*}
 \begin{pmatrix}
 1 & x_1 & x_3 \\ 0 & 1 & x_2 \\ 0 & 0 & 1
 \end{pmatrix}\!,
 \qquad x_i\in\mathbb{R}.
\end{equation*}
An isomorphic representation is obtained by taking new coordinates
\begin{equation*}
 x_1 = x, \qquad x_2 = y, \qquad x_3 = z +\frac12 x_1 x_2,
\end{equation*}
in which the group action is
\begin{equation}
 (x_1,y_1,z_1)\cdot(x_2,y_2,z_2)=\bigg(x_1+x_2,y_1+y_2,z_1+z_2+\frac12(x_1 y_2-x_2 y_1)\bigg).
 \label{groupop}
\end{equation}

This manifold carries a maximally non-integrable distribution spanned
by the vector fields
\begin{equation*}
 X_1 = \partial_x - \frac12 y\partial_z,\qquad
 X_2 = \partial_y + \frac12 x\partial_z,
\end{equation*}
and which define the sub-Riemannian structure. Furthermore, there
exists a sub-elliptic Laplacian $\Delta:=X_1^2+X_2^2$ for which the
fundamental solution to the Poisson equation $\Delta U=\varrho$, with density $\varrho$, is known~\cite{Folland}. This makes the analogy complete because one can
construct the kinetic energy with the sub-Riemannian metric, and take
the potential as the point-source solution $U$;
the Hamiltonian of such Kepler--Heisenberg system is then
\begin{equation}
 \label{KH}
 H = \frac12 \bigg(p_x - \frac12 y p_z\bigg)^2 + \frac12\bigg(p_y+\frac12 x p_z\bigg)^2
 - \frac{\kappa}{\rho}, \qquad
 \rho:=\sqrt{(x^2+y^2)^2+{16}z^2},
\end{equation}
where $\kappa$ is a non-zero real parameter. The properties of this
system have been analysed in \cite{dots, 1,2}, where it was shown that it
has an invariant submanifold given by $z=0$, $p_z=0$,
$p_{\theta} = x p_y -y p_x=0$ on which all trajectories are straight
lines. More importantly, it was also demonstrated that all periodic
solutions must lie on the zero-energy level and that the whole system
is Liouville integrable there. This happens because of the quantity
$J=x p_x+ y p_y +2 z p_z$, for which $\dot{J} = 2H$, and on the level
$H=0$, it is the third integral of motion next to $H$ and
$p_{\theta}$, with which it commutes.
\begin{Remark}\label{rem1}
 The form of potential in~\eqref{KH} is the same as in~\cite{dots} which is the correction of that~in~\cite{1}.
\end{Remark}
We complete the above findings by proving that for all the other
values of energy, the system is not integrable. In addition, the generalisation of the above system to two bodies is straightforward, and we prove its non-integrability as well.


% Original source lines 118--162
\section{The Kepler problem}

In this section we prove that the Kepler problem on the Heisenberg group is not integrable.
This result will follow as a corollary from a more general theorem.
Our proof is based on the Morales--Ramis theorem \cite{Morales:01::b1} which gives
the following necessary conditions for the integrability.
\begin{Theorem} \label{thm:MoRa}
 If a Hamiltonian system is Liouville integrable, with meromorphic
 first integrals, then the identity component of the differential
 Galois group of variational equations along any non-constant
 solution is Abelian.
\end{Theorem}

Two remarks are in order before we jump to application. One is the
``fine print'' of the above theorem: if the variational equation is
not Fuchsian, then only rational first integrals can be~treated. This
will turn out to be the case here, due to the choice of the particular
solution.

Secondly, because the Hamiltonian~\eqref{KH} itself is not
meromorphic, thanks to algebraic potential
\begin{equation*} %\label{potkep}
 V= - \frac{\kappa}{\sqrt{\big(x^2+y^2\big)^2+{16}z^2}},
\end{equation*}
a slight modification to Theorem~\ref{thm:MoRa} is necessary. We describe it in Appendix~\ref{App0}.

When applying the above, the main difficulty is connected
with determination of properties of the differential Galois group of
variational equations. If they can be reduced to a second order equation, then we can use the decisive Kovacic
algorithm~\cite{3}. Sometimes it is possible to~show that the
considered variational equations contain as a subsystem an equation
for which the differential Galois group is known, e.g., hypergeometric
equation and its confluent form. Here~we give a very useful example
which we will apply later.
\begin{Theorem}[H.P.~Rehm, 1979] \label{thm:par}
 Assume that complex parameters $\alpha\neq 0$, $\beta$, and $\gamma$
 of the parabolic cylinder equation
 \begin{equation} \label{eq:5}
 w''(z) - \big(\alpha^2 z^2 +2 \alpha \beta z + \gamma\big) w(z)=0
 \end{equation}
 are such that $\big(\beta^2-\gamma\big)/\alpha$ is not an odd integer. Then
 its differential Galois group is $\mathrm{SL}(2,\mathbb{C})$.
\end{Theorem}
This theorem was proved in \cite{Rehm} and later in~\cite{baby} it was
proved in another way with the help of the Kovacic algorithm.

\setcounter{section}{2}
% Original source lines 341--546
\section{Two-body problem}

Having dealt with the original generalisation of the Kepler problem
proposed in \cite{1}, a natural question arises about the two-body
problem. In the classical Kepler problem, there is no fundamental
difference between one and two bodies: the latter still leads to the
Kepler problem for a single body of reduced mass, revolving around the
center of mass. The reduction is possible due to the symmetries of the Euclidean space, which generate boosts, and correspond closely to the motion: relative positions of two particles follow a geodesic. This is not the case for the Heisenberg group, where the group operation \eqref{groupop} does not preserve geodesics, as discussed in~detail by the authors of~\cite{1}~-- the difference leads them to pose the integrability question also for the two-body case. In what follows, we give a decisive answer: the two-body problem on the Heisenberg group is not integrable.

For two point masses $m_1$ and $m_2$, whose positions are group
elements $g_k = (x_k,y_k,z_k)$ we will take the Hamiltonian to be
\begin{gather}
 H = \frac{1}{2m_1}\bigg(\bigg(p_{x_1}-\frac12 y_1 p_{z_1}\bigg)^2
 + \bigg(p_{y_1}+\frac12 x_1 p_{z_1}\bigg)^2\bigg)\nonumber
 \\ \hphantom{H =}
 {}+ \frac{1}{2m_2}\bigg(\bigg(p_{x_2}-\frac12 y_2 p_{z_2}\bigg)^2 +
 \bigg(p_{y_2}+\frac12 x_2 p_{z_2}\bigg)^2\bigg) +
 V\big(\rho\big(g_1^{-1}\cdot g_2\big)\big),\label{eq:2b}
\end{gather}
where the potential is specified by
\begin{equation*}
 V(\rho) = -\frac{\kappa m_1 m_2}{\rho},\qquad
 \rho(g) = \sqrt{\big(x^2+y^2\big)^2+{16}z^2}.
\end{equation*}

We note that the following first integrals are ``known'':
\begin{alignat*}{3}
& I_1 = p_{x_1} +\frac12 y_1 p_{z_1} + p_{x_2} +\frac12 y_2 p_{z_2},
 \qquad&&
 I_2 = p_{y_1} -\frac12 x_1 p_{z_1} + p_{y_2} -\frac12 x_2 p_{z_2},&
 \\
 &I_3 = p_{z_1}+p_{z_2}=\{I_1,I_2\},
 \qquad&&
 I_4 = y_1 p_{x_1}-x_1 p_{y_1} + y_2 p_{x_2} - x_2 p_{y_2},&
\end{alignat*}
and they satisfy
\begin{equation*}
 \{I_1,I_4\}=I_2, \qquad
 \{I_2,I_4\} = -I_1.
\end{equation*}
Additionally, as for the Kepler--Heisenberg problem,
\begin{equation*}
 J = x_1 p_{x_1} + y_1 p_{y_1} + 2 z_1 p_{z_1} + x_2 p_{x_2} + y_2 p_{y_2} + 2 z_2 p_{z_2},
\end{equation*}
is such that $\dot{J}=2H$. That is, we have 5 first integrals,
although they do not all commute, and on the zero-energy level $J$
becomes the sixth integral. The question, as before, is whether there
exist enough (here: six) commuting integrals.

Now, we make linear canonical transformation
\begin{alignat*}{3}
 &u_1 = \frac{1}{\sqrt{2}}(y_1-\mathrm{i} x_1),\qquad&&
 p_{u_1} = \frac{1}{\sqrt{2}}(p_{y_1}+\mathrm{i}p_{x_1}),&
 \\
& v_1 = \frac{1}{\sqrt{2}}(y_1+\mathrm{i} x_1),\qquad&&
 p_{v_1} =\frac{1}{\sqrt{2}}(p_{y_1}-\mathrm{i}p_{x_1}),&
 \\
& w_1 = z_1 +z_2, \qquad &&
 p_{w_1} = \frac12(p_{z_1}+p_{z_2}),&
 \\
 &u_2 = \frac{1}{\sqrt{2}}(y_2-\mathrm{i} x_2),\qquad&&
 p_{u_2} = \frac{1}{\sqrt{2}}(p_{y_2}+\mathrm{i}p_{x_2}),&
 \\
 &v_2 = \frac{1}{\sqrt{2}}(y_2+\mathrm{i} x_2),\qquad&&
 p_{v_2} =\frac{1}{\sqrt{2}}(p_{y_2}-\mathrm{i}p_{x_2}),&
 \\
& w_2 = z_1 -z_2, \qquad&&
 p_{w_2} = \frac12(p_{z_1}-p_{z_2}).&
\end{alignat*}
In the new variables the Hamiltonian reads
\begin{gather*}
 H = \frac{((p_{w_1}+p_{w_2})u_1-2\mathrm{i}p_{v_1})
 ((p_{w_1}+p_{w_2})v_1+2\mathrm{i}p_{u_1})}{4m_1}
 \\ \hphantom{ H =}
{} + \frac{((p_{w_1}-p_{w_2})u_2-2\mathrm{i}p_{v_2})
 ((p_{w_1}-p_{w_2})v_2+2\mathrm{i}p_{u_2})}{4m_2} -\frac{\kappa m_1 m_2}{\rho_{12}},
\end{gather*}
where
\begin{equation*}
 \rho_{12} = 2\sqrt{(u_1-u_2)^2(v_1-v_2)^2-(v_1 u_2 - v_2 u_1 + 2\mathrm{i}w_2)^2}.
\end{equation*}

The particular solution is almost as before
\begin{equation*}
 p_{w_1}=p_{w_1}(0),\qquad
 p_{w_2} = -2at,\qquad
 a = \frac{m_1 m_2\kappa}{8w_2|w_2|},\qquad
 w_2=w_2(0), \qquad
 w_1=w_1(0),
\end{equation*}
and all other phase variables equal to zero. The solution must not be constant, so $w_2(0)\neq 0$, but other parameters are not restricted.

Linear variations of the variables
$(u_1, p_{v_1},u_2,p_{v_2},v_1,p_{u_1},v_2,p_{u_2},w_1, w_2,
p_{w_1},p_{w_2})$, which we will denote by $\xi=(\xi_1, \ldots, \xi_{12})$, then
satisfy the variational equations
\begin{equation}
 \label{eq:varg}
 \dot \xi = A\xi, \qquad A = \begin{bmatrix}
 A_1 & 0 & 0\\
 0 & A_2 & 0\\
 0 & 0 & A_3
 \end{bmatrix}\!,
\end{equation}
where
\begin{equation*}
 A_1 = \begin{bmatrix}
 \tau -\tau_0& 1 & 0 & 0\\
 (\tau-\tau_0)^2 & \tau-\tau_0 & -1 & 0\\
 0 & 0 & -\mu(\tau+\tau_0) & \mu \\
 1 & 0 & \mu(\tau+\tau_0)^2 & -\mu(\tau+\tau_0)
 \end{bmatrix}\!,\qquad
 A_3 = \begin{bmatrix}
 0 & 0 & 0 & 0\\
 0 & 0 & 0 & 0\\
 0 & 0 & 0 & 0 \\
 0 & 4\mathrm{i}/w_2 & 0 & 0
 \end{bmatrix}\!,
\end{equation*}
and
\begin{equation*}
 A_2 = \begin{bmatrix}
 \tau_0-\tau & 1 & 0 & 0\\
 (\tau_0-\tau)^2 & \tau_0-\tau & 1 & 0\\
 0 & 0 & \mu(\tau+\tau_0) & \mu \\
 -1 & 0 & \mu(\tau+\tau_0)^2 & \mu(\tau+\tau_0)
 \end{bmatrix}\!.
\end{equation*}
To obtain the above form we use the following rescalings
\begin{equation*}
 t = m_1\tau,\qquad
 a = \frac{\mathrm{i}}{m_1},\qquad
 \mu = \frac{m_1}{m_2}, \qquad
 p_{w_1}=2 \mathrm{i} \tau_0.
\end{equation*}

\begin{Theorem} %\label{thm:1}
 If $\mu\neq -1$ then the two-body problem on the Heisenberg group is
 not integrable in the Liouville sense.
\end{Theorem}
\begin{proof}
 If the system generated by~\eqref{eq:2b} is integrable then by
 Theorem~\ref{thm:MoRa}, the identity component of differential
 Galois group of variational equations~\eqref{eq:varg} is Abelian.
 This implies that the same property is shared by the differential Galois groups
 of the subsystems of~\eqref{eq:varg}, which have the form
 $\dot\eta=A_i \eta$, $\eta\in\mathbb{C}^4$, for $i=1,2,3$. We
 consider the first of them. It has particular solution
 \begin{equation}
 \label{par}
 \eta(\tau)=(1, \tau_0-\tau,1, \tau_0+\tau).
 \end{equation}
 Using the d'Alambert method, see~\cite{Walter}, we can reduce the
 dimension of the system by one. But~assuming that $\tau_0=0$ we
 achieve more. Namely, linear transformation $\eta\mapsto Q\eta$ with
 $Q=Q(\tau)$ given by
 \begin{equation*}
 Q = \begin{bmatrix}
 1 & 0 & 0 & 0\\
 -\tau & 1 & 0 & 0\\
 1 & 2\tau & -1 & 0\\
 \tau & -1-2\tau^2 & \tau & -\frac{1}{\mu}
 \end{bmatrix}\!,
 \end{equation*}
 brings it to the form $\dot{\eta}=\widetilde{A}_1\eta$, where
 \begin{equation*}
 \widetilde{A}_1 =Q^{-1}\bigg(A_1 Q -\frac{\mathrm{d}}{\mathrm{d}\tau}Q\bigg) = \begin{bmatrix}
 0 & 1 & 0 & 0\\
 0 & 0 & 1 & 0\\
 0 & p_2 & p_1 & 1\\
 0 & 0 & 0 & 0
 \end{bmatrix}\!,\qquad p_1 = 2(1-\mu)\tau,\qquad p_2 =
 3+\mu+4\mu\tau^2.
 \end{equation*}
 Thus the transformed system has block-triangular structure which is
 quite simple: the first coordinate does not enter, while the fourth
 is constant. In other words, to obtain a particular solution, it is
 enough to assume $\eta_4=0$, and choose the subsystem corresponding to the second and third
 components:
 \begin{gather*}
 \eta_2'(\tau) = \eta_3(\tau),
 \\
 \eta_3'(\tau) = p_2(\tau) \eta_2(\tau) + p_1(\tau)\eta_3(\tau).
 \end{gather*}
 As a single equation it reads
 \begin{equation*}
 \eta_2'' = 2(1-\mu)\tau \eta_2' + \big(3+\mu+4\mu \tau^2\big)\eta_2,
 \end{equation*}
 which, after the change $\eta_2=\exp\big[(1-\mu)\tau^2/2\big]w(\tau)$, becomes
 \begin{equation}
 \label{eq:w2b}
 w''(\tau) -(1+\mu)\big[2 +(1+\mu)\tau^2\big]w(\tau)=0.
 \end{equation}
 It is, again, the parabolic cylinder equation~\eqref{eq:5} with parameters
 \begin{equation*}
 \alpha^2= (1+\mu)^2, \qquad
 \beta=0, \qquad
 \gamma = 2(1+\mu).
 \end{equation*}
 Let us assume that $\mu\neq-1$. Then $\alpha\neq 0$ and
 \[
 \frac{\beta^2-\gamma}{\alpha}=-2\sgn(1+\mu)
 \]
 is not an odd integer. Hence, by Theorem~\ref{thm:par} the
 differential Galois group of equation~\eqref{eq:w2b} is~$\mathrm{SL}(2,\mathbb{C})$. This ends the proof.
\end{proof}


% Original source lines 670--679
\section{Concluding remarks}

Our main goal was to answer the question of Montgomery and Shanbrom about integrability of~the (simple) Kepler problem on the Heisenberg group. The answer turned out to be negative, but several generalisations became immediately apparent. First, the potential had a specific radial/axial symmetry, and a whole general class of such potentials could be included; second, and more important, the two-body problem could be formulated in a natural way. We thus extended the analysis, and managed to show, that with reasonable assumptions those extensions were also non-integrable.

We note that potentials not satisfying condition \eqref{eq:con} can be found, such as
\begin{equation*}
 V = z^2\big(x^2+y^2\big)^2 = z^2\big(\rho^2-16z^2\big), \qquad\text{or}\qquad V=\big(x^2+y^2\big)^2R(z,\rho),
\end{equation*}
where $R(z,\rho)$ is not divisible by $\big(x^2+y^2\big)^2$. Integrability of these potentials remains an open question.
One possible way of investigation of such cases is the application of a variant of the direct method. However, we were unable to find any integrals which were polynomials of low degree in momenta. It remains an open question whether our result can be extended to a wider functional class of first integrals, but each case requires a completely different set of methods than those used here, and as such is a subject for separate investigation.


% Original source lines 682--734
\appendix

\section{Systems with algebraic Hamiltonians}\label{App0}

First, let us remark that Theorem~\ref{thm:MoRa} also holds for a
general Poisson system \cite{akn}. When the Hamiltonian function is algebraic but not meromorphic,
we cannot apply this theorem directly. One solution is to find an extension of the phase space (by including additional variables) in~such a way, that the original Hamiltonian lifts to a meromorphic one, and the extended system is Hamiltonian
with respect to a degenerate Poisson bracket, which reproduces
the original problem.

The construction below is a modification of that given in
\cite{alg}, where the reader will find more details and proofs. Let us
consider an $n$ degrees of freedom Hamiltonian system with canonical
coordinates $q,p\in\mathbb{C}^n$, with algebraic Hamiltonian $H(q,p)$
such that $H(q,p)=K(q,p,u)$, where~$u$ is algebraic over
$\mathbb{C}(q)$ with minimal polynomial $P(u)\in\mathbb{C}(q)[u]$, and
$K(q,p,u)\in \mathbb{C}(q,p,u)$ is a~rational function of its
arguments $x=(q,p,u)\in\mathbb{C}^{2n+1}$. We introduce the following
system
\begin{equation}
 \label{eq:dK}
 \dot x = J(x) \nabla_xK(x),
\end{equation}
where $J(x)$ is $(2n+1)\times(2n+1)$ matrix of the form
\begin{equation*}
 J(x):=\begin{bmatrix}
 0 & \mathbbm{1}_n & 0 \\
 -\mathbbm{1}_n & 0 & \frac{1}{\partial_u P}\nabla_q P \\
 0 & -\frac{1}{\partial_u P} \nabla_q P& 0
 \end{bmatrix}\!,
\end{equation*}
with $\mathbbm{1}_n$ equal to the $n\times n$ identity matrix. It defines the Poisson bracket
\begin{equation*}
 \{f,g\}(x) := (\nabla_xf(x))^{\rm T} J(x)\nabla_xg(x),
\end{equation*}
where $f$ and $g$ are smooth functions. The rank of matrix $J(x)$ is
$2n$ and the only Casimir function of the bracket is $P(u)$.
\begin{Lemma}
{\sloppy If $(q(t),p(t),u(t))$ is a solution of equations~\eqref{eq:dK} with
 $P(u(t))=0$, then $(q(t),p(t))$ is a solution of Hamilton's
 equations
 \begin{equation*}
 %\label{eq:has}
 \dot q= \nabla_q H(q,p), \qquad
 \dot p=-\nabla_p H(q,p).
 \end{equation*}}
\end{Lemma}
We omit the proof, as it is rather direct, and ask the interested reader to follow the explanation and steps given in \cite[Section~2]{alg}. This lemma gives us what is needed, that is we reproduce the original system as a~Hamiltonian one with respect to a degenerate Poisson structure defined
by rational matrix $J(x)$, and with rational Hamiltonian function~$K(x)$.

The above general considerations justify the meromorphic assumptions
of the Morales--Ramis theory, but of course for practical purposes the
calculations can be performed in the original coordinates.
\begin{thebibliography}{99}
\bibitem[1]{akn}
Arnold V.I., Kozlov V.V., Neishtadt A.I., Mathematical aspects of classical and
 celestial mechanics, 3rd ed., \textit{Encyclopaedia of Mathematical Sciences}, Vol.~3, \href{https://doi.org/10.1007/978-3-540-48926-9}{Springer-Verlag}, Berlin, 2006.

\bibitem[2]{borisov:16::}
Borisov A.V., Mamaev I.S., Bizyaev I.A., The spatial problem of 2 bodies on a
 sphere. {R}eduction and stochasticity, \href{https://doi.org/10.1134/S1560354716050075}{\textit{Regul. Chaotic Dyn.}}
 \textbf{21} (2016), 556--580.

\bibitem[4]{dots}
Dods V., Shanbrom C., Self-similarity in the {K}epler--{H}eisenberg problem,
 \href{https://doi.org/10.1007/s00332-021-09709-1}{\textit{J.~Nonlinear Sci.}} \textbf{31} (2021), 49, 15~pages,
 \href{https://arxiv.org/abs/1912.12375}{arXiv:1912.12375}.

\bibitem[5]{baby}
Duval A., Loday-Richaud M., Kova\v{c}i\v{c}'s algorithm and its application to
 some families of special functions, \href{https://doi.org/10.1007/BF01268661}{\textit{Appl. Algebra Engrg. Comm.
 Comput.}} \textbf{3} (1992), 211--246.

\bibitem[6]{Folland}
Folland G.B., A fundamental solution for a subelliptic operator, \href{https://doi.org/10.1090/S0002-9904-1973-13171-4}{\textit{Bull.
 Amer. Math. Soc.}} \textbf{79} (1973), 373--376.

\bibitem[7]{3}
Kovacic J.J., An algorithm for solving second order linear homogeneous
 differential equations, \href{https://doi.org/10.1016/S0747-7171(86)80010-4}{\textit{J.~Symbolic Comput.}} \textbf{2} (1986),
 3--43.

\bibitem[8]{alg}
Maciejewski A.J., Przybylska M., Integrability of {H}amiltonian systems with
 algebraic potentials, \href{https://doi.org/10.1016/j.physleta.2015.08.035}{\textit{Phys. Lett.~A}} \textbf{380} (2016), 76--82.

\bibitem[9]{1}
Montgomery R., Shanbrom C., Keplerian dynamics on the {H}eisenberg group and
 elsewhere, in Geo\-met\-ry, Mechanics, and Dynamics, \textit{Fields Inst.
 Commun.}, Vol.~73, \href{https://doi.org/10.1007/978-1-4939-2441-7_14}{Springer}, New York, 2015, 319--342, \href{https://arxiv.org/abs/1212.2713}{arXiv:1212.2713}.

\bibitem[10]{Morales:01::b1}
Morales-Ruiz J.J., Ramis J.P., Galoisian obstructions to integrability of
 {H}amiltonian systems.~{I}, \href{https://doi.org/10.4310/MAA.2001.v8.n1.a3}{\textit{Methods Appl. Anal.}} \textbf{8} (2001),
 33--95.

\bibitem[11]{Rehm}
Rehm H.P., Galois groups and elementary solutions of some linear differential
 equations, \href{https://doi.org/10.1515/crll.1979.307-308.1}{\textit{J.~Reine Angew. Math.}} \textbf{307--308} (1979), 1--7.

\bibitem[12]{2}
Shanbrom C., Periodic orbits in the {K}epler--{H}eisenberg problem,
 \href{https://doi.org/10.3934/jgm.2014.6.261}{\textit{J.~Geom. Mech.}} \textbf{6} (2014), 261--278, \href{https://arxiv.org/abs/1311.6061}{arXiv:1311.6061}.

\bibitem[14]{Walter}
Walter W., Ordinary differential equations, \textit{Graduate Texts in
 Mathematics}, Vol.~182, \href{https://doi.org/10.1007/978-1-4612-0601-9}{Springer-Verlag}, New York, 1998.

\end{thebibliography}
\end{document}
